\documentclass[12pt]{article}

\usepackage[utf8]{inputenc}
\usepackage[T1]{fontenc}
\usepackage[margin=1in]{geometry}
\usepackage{setspace}
\usepackage{mathptmx} % Times Roman font
\usepackage{booktabs}
\usepackage{array}
\usepackage{hyperref}

\pagestyle{plain}

\begin{document}

\begin{center}
{\Large \textbf{Budget and Budget Justification}}\\[0.5em]
{\normalsize Sterling G. Baird, Department of Mechanical Engineering, Brigham Young University}\\[0.3em]
{\normalsize Grant Period: January 1, 2027 -- August 31, 2029}\\[0.3em]
{\normalsize Total Request: \$110,000}
\end{center}

\bigskip

\section*{Budget Summary}

\begin{center}
\begin{tabular}{l r r r}
\toprule
\textbf{Category} & \textbf{Period 1 (20 mo)} & \textbf{Period 2 (12 mo)} & \textbf{Total} \\
\midrule
\multicolumn{4}{l}{\textit{Student Support (must be $\geq$ 60\% = \$66,000)}} \\
\quad Graduate student stipend + benefits & \$29,989 & \$29,989 & \$59,978 \\
\quad Graduate student tuition (1 student) & \$7,875  & \$8,100  & \$15,975 \\
\quad Undergraduate student support       & \$3,750  & \$2,250  & \$6,000 \\
\textbf{Student support subtotal}         & \textbf{\$41,614} & \textbf{\$40,339} & \textbf{\$81,953} \\
\midrule
PI summer stipend (incl.\ fringe)         & \$8,000  & \$8,000  & \$16,000 \\
Conference travel                         & \$2,000  & \$2,000  & \$4,000 \\
Supplies and materials                    & \$7,529  & \$518    & \$8,047 \\
\midrule
\textbf{Period totals}                     & \textbf{\$59,143} & \textbf{\$50,857} & \textbf{\$110,000} \\
\bottomrule
\end{tabular}
\end{center}

\bigskip

\noindent\textit{Student support: \$81,953 / \$110,000 = 74.5\% ($\geq$ 60\% requirement met)}

\bigskip

\section*{Budget Justification}

\doublespacing

\subsection*{Student Support (\$81,953)}

\textbf{Graduate Student Stipend and Benefits (\$59,978).}
One master's-level graduate student will be supported at \$29,989/year (\$21.36/hr $\times$ 27~hrs/week $\times$ 52~weeks), beginning in Fall~2027. The student will conduct FeS and PbS precipitation experiments in the presence of petroleum-derived compounds (sulfide and organic acids), PHREEQC thermodynamic modeling, and Pitzer parameter regression. The student will gain training in aqueous geochemistry of petroleum-derived species, high-throughput experimentation (Python/Opentrons API), PHREEQC thermodynamic modeling, and nonlinear parameter regression. BYU's M.S.\ program supports thesis completion within 2~years, well-matched to the grant period.

\textbf{Graduate Student Tuition (\$15,975).}
Tuition support for one graduate research assistant covering 30~credits over 2~academic years (15~credits/year): \$7,875 in Period~1 [(\$510 + \$15) $\times$ 15~credits] and \$8,100 in Period~2 [(\$510 + \$30) $\times$ 15~credits], reflecting BYU's per-credit tuition and fee structure for the M.Sc.\ program.

\textbf{Undergraduate Student Support (\$6,000).}
Two undergraduate students will be supported at \$1,500 each per year (\$6,000 total over 32~months, allocated proportionally across budget periods) for part-time research during the academic year and/or summer. Undergraduates will assist with stock solution preparation, image analysis, data entry, and quality-control checks, gaining hands-on research experience with quantitative inorganic chemistry and robotic laboratory automation.

\subsection*{PI Summer Stipend (\$16,000)}

The PI requests 0.5 months of summer salary per year at \$8,000/year including fringe benefits (\$8,000 per grant year, \$16,000 total). This time will be dedicated to mentoring students, analyzing data, and preparing manuscripts. [The \$8,000/year cap including fringe is the maximum allowed by ACS PRF.]

\subsection*{Conference Travel (\$4,000)}

Travel funds of \$2,000 per grant year (\$4,000 total) will support the PI and students presenting research results at national conferences such as the ACS National Meeting. Budget covers registration, transportation, and lodging for one conference per year.

\subsection*{Supplies and Materials (\$8,047)}

Funds will support consumable chemicals (FeSO$_4$, FeCl$_2$, PbCl$_2$, Na$_2$S$\cdot$9H$_2$O, NaHS, NaCl, CaCl$_2$, sodium acetate, sodium propionate, cyclohexane carboxylic acid, pH buffer reagents, ICP-OES standards, ultrapure water) and disposable labware (chemically resistant polypropylene and optically clear polystyrene 96-well microplates, adhesive plate seals, filter-tipped pipette tips for the OT-2). Estimated breakdown: consumable chemicals and reagents (\$2,500), 96-well microplates (\$2,000), filter-tipped pipette tips (\$2,500), and adhesive plate seals, gloves, and miscellaneous consumables (\$1,047). Period~1 requires higher supply costs (\$7,529) for initial chemical inventory procurement, with reduced costs in Period~2 (\$518) for ongoing consumables. Supplies and experimental costs beyond this will be supplemented by the PI's startup equipment and supplies budget, which includes the service fees for ICP-OES usage.

\subsection*{Items Not Included}

No funds are requested for overhead/indirect costs, capital equipment, publication costs, laboratory technicians, contractors, consultants, or visiting faculty, in compliance with ACS PRF policy.

\end{document}
